\section{A uniform limit of finite arithmetic assemblies}
\label{sec:composite-joined-limit}

The completed reference can itself be obtained from finite assembled
experiments. Choose finitely many extension degrees, prepare the corresponding
tagged arithmetic systems, and run the copied primitive protocol in each
degree. Uniform relative randomization supplies a tracial common block on
success. The counting boundary and the degree completion then converge
together in trace norm. This section is a separate draft claim awaiting
its own capped review.

\begin{definition}[Finite degree mixture and success profiles]
Fix an integer base degree $s\geq1$, a real $\beta>1$, and named extensions
$E_d/K$ with $|K|=p^s$ and $[E_d:K]=d$ for $d\geq2$. For an integer
cutoff $D\geq2$ put
\[
 L_D(\beta)=\sum_{d=2}^D d^{-\beta},\qquad
 \pi_d=\frac{d^{-\beta}}{L_D(\beta)}.
\]
Prepare the classical degree tag with these probabilities and run the
copied primitive reference instrument in its selected degree. On success,
uniformly randomize the relative index and explicitly discard that
randomizer label. Every failed degree and stopped primitive/averaging
history remains available.

For $t>1$ define
\[
 w_d(t)=\frac{c_d(t^s)}{dt^{sd}},\qquad
 v_d(t)=\frac{w_d(t)}{s(t-1)},\qquad
 Z_D(\beta,t)=\sum_{d=2}^D d^{-\beta}v_d(t),
\]
where $c_d(x)=\sum_{a\mid d}\mu(d/a)x^a$. Set
$v_d(1)=\varphi(d)/d$ and use that value in $Z_D(\beta,1)$.
On the named Hilbert degree sum, let the common successful density have
blocks
\[
 \rho_{D,t}\big|_{\mathbb C^d}
  =\frac{d^{-\beta}v_d(t)}{dZ_D(\beta,t)}I_d
       \quad(2\leq d\leq D),
\]
and zero elsewhere. Define $L_\infty$, $Z_\infty$ and $\rho_{\infty,t}$
by the corresponding infinite sums when they converge.
\end{definition}
\begin{scope}
The degree prior is additional classical preparation data. These are
common-observable outputs, not identifications of the physical field
multiplicity states. At $t=1$, $\rho_{D,1}$ designates the retained first
grade; it is not normalization of a successful zero-probability event.
Randomization makes this reference stationary. The unrandomized finite
protocol remains the separate Frobenius/coherence witness.
\end{scope}
\provenance{D1621}{definitions.md}{composite joined-limit checker}
{composite-joined-adjudication.md}

The reference instrument used in each degree prepares one period-uniform
field state and two fixed zero ancillas, maps $|x\rangle$ to
$|x,x,x^2\rangle$, and projects onto simultaneous relative-Frobenius
invariants after the primitive-period cut. Its common successful state is
$|0\rangle\langle0|$ and its success probability is $w_d(t)$. Relative
randomization averages its $d$ cyclic translates to $I_d/d$.
For $a_d(t)=c_d(t^s)/t^{sd}$, the three preparation histories within
that degree have probabilities
\[
 1-a_d(t),\qquad (1-1/d)a_d(t),\qquad a_d(t)/d.
\]
They sum to one before multiplication by $\pi_d$. The finite mixture
therefore has total success
\[
 W_D(\beta,t)=\frac{s(t-1)Z_D(\beta,t)}{L_D(\beta)}.
\]
\provenance{D1604, D1607, CMP-BOUNDARY, CMP-TRACE}
{boundary-and-spectrum.md}{composite boundary checker}
{composite-boundary-adjudication.md}

\begin{theorem}[A joined counting and degree limit of assembled quantum states]
\statusSketch\quad
For every integer $s\geq1$ and real $\beta>1$, the finite degree-mixture
protocol is normalized with its stated histories. Its rescaled success
factors satisfy $0<v_d(t)\leq1$ for all $d\geq2,t>1$ and extend
continuously to $v_d(1)=\varphi(d)/d$.

For every fixed $T>1$, the infinite normalizer converges uniformly on
$1\leq t\leq T$. Its density satisfies
\[
 \rho_{\infty,t}\longrightarrow\rho_\beta
       \quad\text{in trace norm as }t\downarrow1,
\]
where
\[
 \rho_\beta\big|_{\mathbb C^d}
 =\frac{\varphi(d)}{Z_{\rm deg}(\beta)d^{\beta+2}}I_d,\qquad
 Z_{\rm deg}(\beta)=\frac{\zeta(\beta)}{\zeta(\beta+1)}-1.
\]
For every integer $D\geq2$ and $1\leq t\leq T$,
\[
 \|\rho_{D,t}-\rho_{\infty,t}\|_1
 \leq \frac{2^{\beta+2}T^s}{\beta-1}\,D^{1-\beta}.
\]
Consequently both iterated limits and every path $D\to\infty$, $t\downarrow1$
give the same $\rho_\beta$. With $h=\log t$, their first-order success rate
tends to
\[
 \frac{sZ_{\rm deg}(\beta)}{\zeta(\beta)-1}.
\]
Probabilities for each fixed bounded common effect converge as well,
including after any fixed number of relative Frobenius iterations.
\end{theorem}
\begin{scope}
The degree prior and common-observable comparison are specified choices.
The statement includes no arbitrary characteristic-dependent physical
continuation and requires no total endpoint tensor or direct-sum operation.
The raw event at $t=1$ still has probability zero. The cutoff bound is a
uniform estimate and need not be sharp or smaller than the trivial bound two.
\end{scope}
\provenance{CMP-JOINT-LIMIT}{joined-limit.md, Sections 1--4}
{composite joined-limit checker}{composite-joined-adjudication.md}

\begin{proof}[The uniform estimate]
The positive period polynomials satisfy
$c_d(x)\leq x^d-1$ for $d>1,x\geq1$, since their divisor sum is $x^d$
and the degree-one term is $x$. Also
\[
 1-t^{-sd}=(t-1)\sum_{j=1}^{sd}t^{-j}\leq sd(t-1).
\]
Together these give $v_d(t)\leq1$. Dividing the polynomial $c_d(t^s)$
by $t-1$ gives its continuous endpoint value $\varphi(d)/d$.
Thus the summands of $Z_\infty$ are dominated uniformly by $d^{-\beta}$.

The degree-two block bounds every normalizer away from zero on the compact
interval:
\[
 v_2(t)=\frac1{2s}\sum_{j=1}^s t^{-j}\geq\frac12T^{-s},\qquad
 Z_\infty(\beta,t)\geq2^{-\beta-1}T^{-s}.
\]
Trace norm between these scalar-block densities is the sum of the absolute
differences of their block probabilities. Removing and renormalizing the
tail gives exactly
\[
 \|\rho_{D,t}-\rho_{\infty,t}\|_1
 =2\frac{Z_\infty(\beta,t)-Z_D(\beta,t)}{Z_\infty(\beta,t)}.
\]
The integer-cutoff bound for $\sum_{d>D}d^{-\beta}$ proves the displayed
uniform estimate. A finite head is continuous at $t=1$, and the tail is
uniformly small, giving the remaining trace-norm limit. Finally,
$(t-1)/\log t\to1$ gives the retained success rate. Unitary Frobenius
preserves trace norm, so fixed-effect convergence follows.
\end{proof}

This construction gives the completed state as the limit of actual finite
assemblies performed before specialization. Its common randomized state
is stationary under Frobenius, while the quantum automorphism remains
nontrivial on the available finite-block preparations and observables.
The representation has supplied an infinite quantum system without
requiring a field structure or a conventional arithmetic composition law
at characteristic one.
